\chapter{Kommentierte Quellen}

Achtundzwanzig Werke, eingeordnet: Zu jedem steht, wofür man es aufschlägt und was es leistet, das andere nicht leisten.

\lit{1}{Abu-Mostafa, Y., Magdon-Ismail, M., Lin, H.-T.: \textit{Learning From Data}. AMLBook 2012.}{Kurz, klar und auf die Kernfrage dieses Bandes ausgerichtet: Wann ist Lernen überhaupt möglich? Die VC-Theorie wird hier so zugänglich erklärt wie nirgends sonst. Etwa zweihundert Seiten; die beste Brücke zwischen Teil~I und Teil~III.}

\lit{2}{Belkin, M., Hsu, D., Ma, S., Mandal, S.: \textit{Reconciling modern machine-learning practice and the classical bias--variance trade-off}. PNAS 116, 2019.}{Der Aufsatz, der den doppelten Abstieg benannte. Zeigt die Kurve aus Kapitel~14 an linearen Modellen, Zufallsmerkmalen und Netzen. Kurz und für Nichtfachleute lesbar.}

\lit{3}{Bishop, C.: \textit{Pattern Recognition and Machine Learning}. Springer 2006.}{Das Lehrbuch aus wahrscheinlichkeitstheoretischer Sicht. Mischmodelle und EM (Kapitel~9), Kerne und Gaußprozesse sind besonders gut dargestellt. Anspruchsvoll, mit vielen Abbildungen; die Fundstelle für die achte Auslassung.}

\lit{4}{Breiman, L.: \textit{Statistical Modeling: The Two Cultures}. Statistical Science 16, 2001.}{\textbf{Der Streitaufsatz, der das Verhältnis von Statistik und Lernen neu bestimmte.} Breiman wirft der Statistik vor, an Modellen festzuhalten, wo Vorhersagegüte zählt. Die angehängten Erwiderungen --- unter anderem von Cox und Efron --- sind ebenso lesenswert wie der Aufsatz selbst.}

\lit{5}{Christodoulou, E. u.\,a.: \textit{A systematic review shows no performance benefit of machine learning over logistic regression for clinical prediction models}. J.\ Clin.\ Epidemiol.\ 110, 2019.}{Die Übersicht über $71$ Vergleichsstudien aus Kapitel~5 und 15. Unverzichtbar für jeden, der ein lernendes Verfahren in der Klinik einsetzen will: Sie zeigt, wie oft der scheinbare Vorteil aus der Prüfung stammt.}

\lit{6}{Cover, T., Hart, P.: \textit{Nearest neighbor pattern classification}. IEEE Trans.\ Inf.\ Theory 13, 1967.}{Die Originalarbeit zu Satz 10.2, sieben Seiten. Enthält den Grenzübergang, den Kapitel~10 ausließ, und die Formulierung, dass die Hälfte der Information im nächsten Nachbarn steckt.}

\lit{7}{Deisenroth, M., Faisal, A., Ong, C.: \textit{Mathematics for Machine Learning}. Cambridge 2020.}{Die mathematischen Werkzeuge des Faches --- lineare Algebra, Analysis, Wahrscheinlichkeit, Optimierung --- auf das Lernen zugeschnitten. Für Leser, die Band 2 und 19 nicht zur Hand haben. Frei verfügbar; die Fundstelle für Eckart-Young.}

\lit{8}{Ertel, W.: \textit{Grundkurs Künstliche Intelligenz}. 5.\ Aufl., Springer Vieweg 2021.}{\textbf{Auf Deutsch, als Überblick.} Behandelt maschinelles Lernen im Rahmen der künstlichen Intelligenz insgesamt, mit Logik, Suche und Schließen. Gut lesbar, mit vielen Übungen; weniger Tiefe in der Theorie als die englischen Lehrbücher.}

\lit{9}{Frochte, J.: \textit{Maschinelles Lernen --- Grundlagen und Algorithmen in Python}. 3.\ Aufl., Hanser 2020.}{\textbf{Auf Deutsch und praktisch.} Die Verfahren dieses Bandes mit lauffähigem Code, von der Regression bis zu Netzen. Für Leser, die das Gelesene selbst ausprobieren wollen.}

\lit{10}{Goodfellow, I., Bengio, Y., Courville, A.: \textit{Deep Learning}. MIT Press 2016.}{Das Standardwerk zu tiefen Netzen, von dreien, die das Gebiet geprägt haben. Rückpropagation, Faltungsnetze, Regularisierung, Optimierung. Teilweise vom Stand überholt, in den Grundlagen weiterhin maßgeblich. Frei verfügbar.}

\lit{11}{Hastie, T., Tibshirani, R., Friedman, J.: \textit{The Elements of Statistical Learning}. 2.\ Aufl., Springer 2009.}{\textbf{Das Nachschlagewerk des statistischen Lernens.} Enthält fast jedes Verfahren dieses Bandes mit Herleitung und Diskussion, geschrieben von den Erfindern von Lasso, Boosting-Varianten und Bäumen. Anspruchsvoll; frei verfügbar.}

\lit{12}{James, G., Witten, D., Hastie, T., Tibshirani, R.: \textit{An Introduction to Statistical Learning}. 2.\ Aufl., Springer 2021.}{\textbf{Die sanftere Fassung von [11]:} dieselben Verfasser, weniger Mathematik, mehr Anschauung. Zum Einstieg das beste Buch des ganzen Faches. Frei verfügbar, mit Übungen in R und Python.}

\lit{13}{Kapoor, S., Narayanan, A.: \textit{Leakage and the reproducibility crisis in machine-learning-based science}. Patterns 4, 2023.}{Die Sammlung von Datenlecks aus siebzehn Fachgebieten, auf die Kapitel~3 verweist. Mit einer Einteilung der Leckarten und einer Prüfliste. \textbf{Wer lernende Verfahren in der Forschung einsetzt, sollte diese Liste abarbeiten.}}

\lit{14}{Mohri, M., Rostamizadeh, A., Talwalkar, A.: \textit{Foundations of Machine Learning}. 2.\ Aufl., MIT Press 2018.}{Die Lerntheorie mit vollständigen Beweisen --- Rademacher-Komplexität, Randschranken, Boosting. Die Vertiefung zu Teil~III für Leser mit mathematischer Ausbildung.}

\lit{15}{Murphy, K.: \textit{Probabilistic Machine Learning: An Introduction}. MIT Press 2022.}{Das umfassendste aktuelle Lehrbuch, aus wahrscheinlichkeitstheoretischer Sicht. Reicht von der Regression bis zu Sprachmodellen. Umfangreich und aktuell; frei verfügbar.}

\lit{16}{Obermeyer, Z., Powers, B., Vogeli, C., Mullainathan, S.: \textit{Dissecting racial bias in an algorithm used to manage the health of populations}. Science 366, 2019.}{Der Fall aus Kapitel~15. \textbf{Ein Musterbeispiel dafür, wie eine falsch gewählte Zielgröße zu systematischer Benachteiligung führt} --- ohne jeden Fehler im Verfahren. Kurz und mit bemerkenswert klarer Methodik.}

\lit{17}{Prince, S.: \textit{Understanding Deep Learning}. MIT Press 2023.}{Die aktuellste und anschaulichste Darstellung tiefer Netze, mit ausgezeichneten Abbildungen. Behandelt den doppelten Abstieg, implizite Regularisierung und die Frage, warum Netze verallgemeinern --- die Vertiefung zu Kapitel~14. Frei verfügbar.}

\lit{18}{Rasmussen, C., Williams, C.: \textit{Gaussian Processes for Machine Learning}. MIT Press 2006.}{Das Werk zu Gaußprozessen aus dem Forschungskasten von Kapitel~6: Kerne mit Unsicherheit. Für Anwendungen, in denen man wissen muss, wo das Verfahren nichts weiß. Frei verfügbar.}

\lit{19}{Richter, S.: \textit{Statistisches und maschinelles Lernen}. Springer Spektrum 2019.}{\textbf{Auf Deutsch und mathematisch.} Aus einer Vorlesung hervorgegangen, mit Beweisen und durchgerechneten Beispielen. Die deutschsprachige Entsprechung zu [11] in kompakterer Form.}

\lit{20}{Riley, R. u.\,a.: \textit{Calculating the sample size required for developing a clinical prediction model}. BMJ 368, 2020.}{Die Formeln für die Mindestgröße klinischer Vorhersagestudien aus Kapitel~2. Sie ersetzen die Zehnerregel durch eine Rechnung, die den Optimismus ausdrücklich begrenzt. \textbf{Vor jeder Studie zu einem Vorhersagemodell zu lesen.}}

\lit{21}{Russell, S., Norvig, P.: \textit{Künstliche Intelligenz --- Ein moderner Ansatz}. 4.\ Aufl., Pearson 2023.}{\textbf{Die deutsche Ausgabe des weltweit verbreitetsten Lehrbuchs der künstlichen Intelligenz.} Maschinelles Lernen ist dort ein Teil unter vielen; wer den größeren Zusammenhang sucht --- Planen, Schließen, Handeln unter Unsicherheit ---, findet ihn hier.}

\lit{22}{Schölkopf, B., Smola, A.: \textit{Learning with Kernels}. MIT Press 2002.}{Das Standardwerk zu Kernmethoden, von zwei der Hauptentwickler. Enthält den Satz von Mercer mit Beweis, die Theorie der Stützvektormaschinen und die Kern-Hauptkomponenten. Die Vertiefung zu Kapitel~6.}

\lit{23}{Seymour, C. u.\,a.: \textit{Derivation, validation, and potential treatment implications of novel clinical phenotypes for sepsis}. JAMA 321, 2019.}{Die Studie mit den vier Sepsis-Phänotypen aus Kapitel~12. Ein gutes Beispiel für Gruppierung mit anschließender Prüfung an Größen, die nicht in die Gruppierung eingingen --- und für die offenen Fragen, die danach bleiben.}

\lit{24}{Shalev-Shwartz, S., Ben-David, S.: \textit{Understanding Machine Learning: From Theory to Algorithms}. Cambridge 2014.}{\textbf{Für die Beweise die maßgebliche Quelle.} Enthält alles, was Teil~III ausließ --- Hoeffding, Sauer, den Fundamentalsatz mit Symmetrisierung ---, in klarer, schrittweiser Darstellung. Frei verfügbar.}

\lit{25}{Vapnik, V.: \textit{The Nature of Statistical Learning Theory}. 2.\ Aufl., Springer 2000.}{Die Theorie aus der Sicht ihres Begründers, knapp und ohne Beweise, mit philosophischen Einschüben. Eigenwillig, aber aufschlussreich: Man versteht, warum Vapnik den Rand für das eigentliche Prinzip hält.}

\lit{26}{Wong, A. u.\,a.: \textit{External validation of a widely implemented proprietary sepsis prediction model in hospitalized patients}. JAMA Intern.\ Med.\ 181, 2021.}{\textbf{Der Leitfall dieses Bandes.} Die Nachprüfung eines Sepsis-Warnsystems an rund $38\,000$ Aufenthalten. Kurz, sachlich --- und eine Warnung für jeden, der ein kommerzielles Vorhersagesystem ungeprüft übernimmt.}

\lit{27}{Zech, J. u.\,a.: \textit{Variable generalization performance of a deep learning model to detect pneumonia in chest radiographs}. PLOS Medicine 15, 2018.}{Das Beispiel für Abkürzungen aus Kapitel~14: Ein Netz erkannte die Klinik statt der Krankheit. Zeigt zugleich, wie man das mit einfachen Mitteln aufdeckt.}

\lit{28}{Zhang, C., Bengio, S., Hardt, M., Recht, B., Vinyals, O.: \textit{Understanding deep learning (still) requires rethinking generalization}. Comm.\ ACM 64, 2021.}{Die Zeitschriftenfassung des Versuchs mit vertauschten Etiketten aus Kapitel~14, mit einem Rückblick auf die Debatte, die er auslöste. Zehn Seiten, die zeigen, wo die Theorie heute steht.}

\begin{merksatz}
\boxtitel{Wenn Sie nur drei Bücher lesen}
\begin{center}
\renewcommand{\arraystretch}{1.6}
\begin{tabular}{@{}P{1.0cm}P{3.4cm}P{3.8cm}@{}}
\toprule
\textbf{[12]} & \textbf{James u.\,a.} & Für den Einstieg\\
\textbf{[24]} & \textbf{Shalev-Shwartz} & \textbf{Für die Beweise}\\
\textbf{[13]} & \textbf{Kapoor, Narayanan} & \textbf{Für die Wirklichkeit}\\
\bottomrule
\end{tabular}
\end{center}

\begin{center}
\begin{minipage}{0.88\textwidth}
\kursiv{James und Mitarbeiter erklären die Verfahren, Shalev-Shwartz und Ben-David beweisen, warum sie funktionieren. Kapoor und Narayanan zeigen, wie oft sie in der Praxis nicht funktionieren --- \textbf{und die dritte Lektüre ist für die Medizin die wichtigste.} Alle drei sind frei verfügbar.}
\end{minipage}
\end{center}

\begin{center}
\begin{minipage}{0.88\textwidth}
\kursiv{Nach Interessenlage: [11] Hastie u.\,a. als Nachschlagewerk; [17] Prince für tiefe Netze; [20] Riley vor jeder eigenen Studie; [26] Wong als Mahnung. Auf Deutsch: [19] Richter für die Mathematik, [9] Frochte für die Praxis, [8] Ertel und [21] Russell und Norvig für den Überblick.}
\end{minipage}
\end{center}
\end{merksatz}

\bigskip
\begin{flushright}
\itshape\color{grau}
Die Daten sagen, was war.\\
Das Modell sagt, was sein könnte.\\
Nur die Prüfung sagt,\\
ob es für den nächsten Fall gilt.
\end{flushright}
